\documentclass{article}
\usepackage{amsmath}
\usepackage[hidelinks]{hyperref}
\title{A local observation can arrive after the last useful steering cut}
\date{}
\begin{document}
\maketitle

\paragraph{Same physical plant, two causal worlds.}
Take the actual position and velocity at the first request in the
published planar \href{https://arxiv.org/html/2609.29939v1}{KTH
managed-departure simulator}. The counterfactual removes the retiring
robot and parks its physical body. The singleton survivor and the
original retiring robot keep their same source states. In two
author-selected branches, the hidden three-robot component is
translated $\pm0.8$ in the direction perpendicular to the initial
nearest cross pair and keeps its source velocities. This is not
an observed pair from the original published run. To show common
initial history is feasible, over three pre-cut units the remote
component drifts with its inherited velocity and adds the same
smooth $3z^2-2z^3$ shift to each member, starting with zero added
position and ending with zero added velocity. Both histories have
the same earlier joint position and velocity at the start of that
interval. An unchanged directly sensed retiring robot does not
forward those remote changes. This no-forwarding premise is about
the admitted interface, not a claim that the original controller
would fail to forward. The certified pre-cut local-to-remote
distance is at least $1.9964>1.6$, the old bridge's named edges
stay in their ordinary range, and all designed motion respects
the same speed and acceleration norm caps of one.

At the actual fault cut, each surviving robot senses positions
perfectly and continuously within nominal range $S=1.6$. There is
no other report or side channel revealing which pre-cut branch
occurred. The past direction of the failed robot is equal in the
two branches. Preserve the continuing internal edges and
clearance $0.45$ from other robots and the failed body. The
terminal requirement is a new cross link at margin range
$R=1.54$, with all survivors at rest at $T=1.11$. This is a
geometric settled bridge, not packet delivery or field sensing.
Use ordinary pre-coordination or common random seeds if desired;
neither changes the information available before contact.

\paragraph{Departure from full-information feasibility.}
The inherited warm-state constructive law brakes, then translates
each component toward a closest cross pair under the shared caps.
With the sign known this produces a verified trajectory in each
branch, reaches the required link in about $1.108498$, and then
holds it. All pairs are conservatively checked over the continuous
motion; clearance to any threatened body is at least $0.7240$.
Holding after common braking while a reliable branch distinction
arrives at time $0.002$ still gives a witness finishing about
$1.1095433$. The alternative medium and shared conditional
codebook are stipulated; no mechanism to deliver that bit in
time is established. No added bit is needed once all branch
positions are supplied as full-state input.

\paragraph{Why even a perfect passive read arrives too late.}
For any trajectory with speed capped by $V$ and acceleration
bounded by $A$ that stops at time $T$, speed up to any cut $t$
is bounded by $\min\{V,\|v(0)\|+At,A(T-t)\}$. Integrating the
bound from zero to $t$ gives a conservative maximum distance
traveled. On either branch, even the sum of this allowance for
every cross pair falls more than $0.0043$ short of closing to
sensing range by $\tau=0.85$. Any shared policy which succeeds
in both branches must therefore have the same singleton
trajectory and velocity through $\tau$. No blackout is imposed
after it is physically possible to sense.

Write $x_-,x_+$ for singleton terminal positions along two
hypothetical successful runs. Velocity difference is zero at the
common cut and at terminal rest. At time $t$ after that common
cut its norm is bounded by
$2A\min\{t,T-\tau-t\}$. Integration gives
\[
\|x_--x_+\|\leq A(T-\tau)^2/2=0.0338.
\]
At this deadline only the nearest-index right robot could be a
valid partner in either branch, even granting arbitrary motion to
the other robots. Its actual terminal location is inside the ball
around its own starting position with radius given by its
whole-horizon travel bound. Combining this with terminal range
$R$ places $x_-$ and $x_+$ respectively in two candidate disks
centered at the two branch positions of that right robot.
Both are also in the common singleton reachable disk. Consequently
the midpoint $m=(x_-+x_+)/2$ lies in the singleton disk and must
belong to both right disks enlarged by $0.0169$.

For any point in two disks with centers $b,c$ and radii $r,s$,
the averaged squared-distance identity bounds its distance to
$(b+c)/2$ by the square root of
$(r^2+s^2)/2-\|b-c\|^2/4$. This enlarged enclosure is still at
least $0.0060$ away from the singleton disk. The midpoint
therefore cannot exist. An independent explicit two-dimensional
circle-arrangement intersection check agrees. Lean checks the
identity and enclosure. Numerical calculations use slack on
the pinned floating inputs, not bit-exact physical coordinates.
Because the necessary terminal sets are incompatible even after
granting motion based on arbitrary information after $\tau$,
not even a stronger active planner using the same range and
prior information can guarantee both outcomes. A randomized
policy that guarantees both with probability one would induce
a shared seed that succeeds on both finite branches, also
contradicting the proof.

\paragraph{Relation to the intended crash baseline.}
The \href{https://en.endo-lab.jp/research/detail/499}{authors'
description of Maeda, Endo and Matsuno's recovery}
uses local neighboring motion, the failed member's preceding
direction, and optional waiting. The constructed pair holds
that cue fixed, so any policy with just the admitted local
signals is subject to the same lower bound. This does not claim
a replay of those authors' full controller, a failure of their
actual experiments or the same terminal deadline in their
paper. The result separates loss of physical feasibility from
late information using a sustained bounded-acceleration test.
The two selected branches and their observation model are
development cases, not an independent mission or an R6
acceptance decision.
\end{document}
